\chapter*{Заключение}                       % Заголовок
\addcontentsline{toc}{chapter}{Заключение}  % Добавляем его в оглавление

Диссертационная работа посвящена разработке и исследованию методов безопасного обучения с подкреплением для задач, в которых агент должен одновременно исследовать среду, достигать целевой награды и соблюдать ограничения безопасности. Основное внимание уделено трём направлениям: использованию моделей мира для предварительной оценки последствий действий, иерархической декомпозиции поведения для задач с большим горизонтом и переходу к мультимодальной постановке, где цели и ограничения задаются естественным языком.

В работе проведён обзор современных методов безопасного обучения с подкреплением и показано, что существующие подходы сталкиваются с рядом ограничений. Неиерархические безопасные стратегии часто становятся чрезмерно консервативными и плохо решают задачи с большим горизонтом. Методы на основе моделей мира позволяют заранее оценивать последствия действий, но чувствительны к ошибкам модели и выбору порогов безопасности. Методы, использующие априорные знания о среде, могут обеспечивать более строгий контроль безопасности, однако требуют известных динамических моделей, явных функций ограничений или дополнительной информации о среде, что ограничивает их применимость в реальных робототехнических задачах.

В рамках направления безопасного обучения с подкреплением на основе модели мира исследован подход SPOWL, в котором модель мира используется для оценки награды, стоимости и долгосрочного риска кандидатных действий. Показано, что безопасное исследование должно учитывать не только мгновенную стоимость, но и оценки долгосрочной стоимости, позволяющие оценивать риск за пределами короткого горизонта планирования. Экспериментальный анализ показал, что способ агрегации ансамбля функций оценки долгосрочной стоимости существенно влияет на баланс между безопасностью и эффективностью: минимальная агрегация недооценивает риск, максимальная приводит к избыточной консервативности, а усреднение обеспечивает наиболее устойчивое поведение агента.

Для решения проблемы безопасного исследования в задачах с большим горизонтом разработан метод SG-Safe. Метод использует стратегию генерации промежуточных подцелей, которая применяется на этапе обучения и направляет безопасную стратегию к более информативным областям пространства состояний. При этом на этапе исполнения используется только итоговая безопасная стратегия, поэтому метод не требует онлайн-планирования и дополнительных вычислительных затрат. Эксперименты на Safety Gym и POLAMP показали, что SG-Safe повышает успешность решения навигационных задач при сохранении низкого уровня нарушений безопасности. В среднем на Safety Gym метод достигает доли успешных эпизодов $0.87$ при частоте столкновений $0.03$, а на POLAMP демонстрирует преимущество в задачах с длинными маршрутами и узкими коридорами.

Дальнейшее развитие иерархического подхода реализовано в методе ITES. В отличие от SG-Safe, метод ITES вводит явный контроль безопасности как на уровне генерации подцелей, так и на уровне их достижения. Верхнеуровневая стратегия генерирует промежуточные подцели с учётом модели стоимости и высокоуровневого критика стоимости, а низкоуровневая стратегия использует модель мира для проверки безопасности коротких воображаемых траекторий. Эксперименты на SafeAntMaze и SafePusher показали, что ITES превосходит неиерархические методы безопасного обучения с подкреплением в задачах с большим горизонтом, где многие методы сравнения либо нарушают ограничения, либо не достигают цели. В SafeAntMaze C-shape метод достигает доли успешных эпизодов $0.88$ при стоимости $15.2$, а в SafePusher~--- доли успешных эпизодов $0.42$ при стоимости $12.8$, сохраняя ограничения безопасности.

Отдельно в диссертации рассмотрена задача безопасного обучения с подкреплением с естественно-языковыми целями и ограничениями. Для этой постановки разработан бенчмарк Caged CrafText, расширяющий CrafText до задач, где агент получает не только инструкцию цели, но и текстовое ограничение безопасности. Бенчмарк включает бюджетные, последовательные, математические и отношенческие ограничения, проверяющие память, безопасное исследование, рассуждение, привязку инструкции к объектам и состояниям среды, а также долгосрочное планирование. Эксперименты с PPO, PPO-LAG, PPO-Cost и агентом на основе большой языковой модели показали, что существующие подходы не решают задачу устойчивого совмещения следования инструкциям и безопасности. PPO достигает награды, но систематически нарушает ограничения; PPO-LAG снижает стоимость, но часто становится чрезмерно консервативным; агент на основе большой языковой модели лучше интерпретирует инструкции, но сохраняет высокий уровень нарушений из-за ошибок привязки текстового ограничения к состояниям и действиям среды.

Полученные результаты показывают, что безопасное обучение с подкреплением требует не только оптимизации численной функции стоимости, но и более структурированных механизмов принятия решений. Модели мира позволяют заранее оценивать риск действий, иерархическая структура улучшает исследование в задачах с большим горизонтом, а мультимодальные бенчмарки выявляют ограничения современных агентов при работе с естественно-языковыми требованиями безопасности.

Несмотря на полученные результаты, остаются открытыми несколько направлений дальнейшей работы. Во-первых, требуется дальнейшее исследование устойчивости моделей мира при увеличении горизонта планирования и при переносе методов из симуляторов на реальные робототехнические платформы. Во-вторых, перспективным направлением является разработка методов, которые объединяют иерархическую декомпозицию, безопасное исследование и перепланирование во время исполнения при изменении условий среды. В-третьих, для задач с естественно-языковыми ограничениями необходимы агенты, способные надёжно связывать текстовые инструкции с объектами, действиями и динамикой среды, а также уточнять неоднозначные ограничения через взаимодействие с человеком.

Таким образом, диссертационная работа развивает методы безопасного обучения с подкреплением от классической постановки с численно заданными ограничениями к более общей постановке, включающей модели мира, иерархическую структуру поведения и естественно-языковые ограничения. Полученные результаты могут быть использованы при разработке автономных робототехнических систем, которым необходимо адаптироваться к новым условиям, исследовать среду и соблюдать заданный уровень безопасности.
